\section{Brown--Peterson coordinates and the Adams--Novikov resolution}
\label{fw:section}

Throughout the proof,
\begin{equation}
 p=1009,\qquad q=2p-2=2016,\qquad
 u_i=\frac{p^i-1}{p-1}\quad(i\geq 0).
 \label{fw:parameters}
\end{equation}
The integer $1009$ is prime: its square root is less than $32$, and none of
$2,3,5,7,11,13,17,19,23,29,31$ divides it.  Put
\[
 R=BP_* =\mathbb Z_{(p)}[v_1,v_2,\ldots],\qquad
 \Gamma=BP_*BP=R[t_1,t_2,\ldots],\qquad
 I_m=(p,v_1,\ldots,v_{m-1}).
\]
Both $v_i$ and $t_i$ have internal degree $qu_i$.  We call internal degree
divided by $q$ the \emph{weight}.  A cobar class of cohomological degree $s$
and weight $w$ has stem $qw-s$.

We use the ordinary Brown--Peterson spectrum, its polynomial cooperation
algebra, and a strictly unital associative model of its multiplication.
For the identification of these polynomial rings with the homotopy and
cooperations of $BP$, see
\citet[Theorems~4.1.18--4.1.19]{Ravenel2004}. The Hazewinkel recursion
and the integrality of its polynomial generators are given in
\citet[Equation~A2.2.1 and Theorem~A2.2.3]{Ravenel2004}; the formal-sum
coordinate convention is explained in the proof of Theorem~A2.1.27 there.
The formal-group foundations are due to
\citet[Section~5, Theorem~4]{Quillen1969}; the explicit generators and
cooperations are developed in \citet[Sections~3.1 and~4.4]{HazewinkelIII}.
A coherent associative model exists, in particular, by forgetting structure
from \citet[Theorem~1.1]{BasterraMandell2013}.  No commutative or
$E_\infty$ model of $BP$ is required below.  We give the particular coordinate,
convergence, and pairing arguments needed for the construction.

\subsection{Coordinates and additive ranges}

Let $a_i$ be the coefficients in the formal-group sum expression
\[
 [p]_F(x)=\mathop{\sum}\nolimits_F a_i x^{p^i},\qquad a_0=p.
\]
The strict coordinate change is written in the same convention as
$\sum_F t_i x^{p^i}$, with $t_0=1$.  Reversing the convention for composition
reverses the cobar conventions without changing any of the arguments.

\begin{lemma}[Leading coefficients and low cooperations]
\label{fw:coordinates}
For each $m\geq 1$, $I_m$ is an invariant ideal and $v_m$ is primitive
modulo $I_m$.  Modulo $I_m$, the coproducts of $t_i$, $i\leq m$, are the
additive-composition coproducts.  At the specialization that kills $p$ and
all $v_i$, the resulting Hopf algebra is
\begin{equation}
 P=\mathbb F_p[t_1,t_2,\ldots],\qquad
 \Delta(t_i)=\sum_{\ell=0}^i t_\ell\otimes t_{i-\ell}^{p^\ell}.
 \label{fw:additive-hopf}
\end{equation}
\end{lemma}

\begin{proof}
Write the logarithm as $\log_F(x)=\sum_{i\geq0}l_i x^{p^i}$, where $l_0=1$.
The Hazewinkel recursion is
\[
 pl_n=v_n+\sum_{i=1}^{n-1}l_i v_{n-i}^{p^i}.
\]
Applying the logarithm to the formal sum for $[p]_F$ and comparing the
coefficient of $x^{p^n}$ shows that the coefficient of the indecomposable
$v_n$ in $a_n$ is
\begin{equation}
 c_n=1-p^{p^n-1}.
 \label{fw:coordinate-unit}
\end{equation}
Indeed the two terms involving $l_n$ are $pl_n$ and $p^{p^n}l_n$, and the
coefficient of $v_n$ in $l_n$ is $1/p$.  All other terms involve lower
$v_i$.  The coefficients $a_n$ are integral, so the grading and polynomial
independence of the $v_i$ give
$a_n-c_nv_n\in(v_1,\ldots,v_{n-1})R$.  Each $c_n$ is a $p$-local unit.

The invariance and primitivity follow inductively.  The element $p$ is
invariant.  Once the lower coefficients vanish, the first surviving term
of the $p$-series is $c_m v_m x^{p^m}$.  A strict coordinate change has
leading term $x$, so comparison of this first surviving term on the two
sides of the coordinate-change identity fixes $c_m v_m$, and hence $v_m$.
This proves both invariance of the next ideal and the asserted primitivity.

After all base generators have been killed, the formal group is additive,
and composition of additive $p$-polynomials gives
Equation~\eqref{fw:additive-hopf}.  Modulo $I_m$, a correction to this
coproduct must have a positive base coefficient, of weight at least $u_m$.
For $i<m$ its weight is already too large.  For $i=m$ the only possibility
is a constant multiple of $v_m$ with no positive-weight cooperation factor.
Applying either counit excludes that term.  This proves the stated range.
\end{proof}

\begin{lemma}[The relation supplied by $v_5$]
\label{fw:v5-relation}
In the cobar cohomology with coefficients $R/I_4$, there is a scalar
$e\in\mathbb F_p$ such that
\begin{equation}
 v_4\bigl(h_{1,4}+e v_4^{p-1}h_{1,0}\bigr)=0,
 \label{fw:v5-cohomology-relation}
\end{equation}
where $h_{1,j}$ is represented by the primitive $t_1^{p^j}$.
\end{lemma}

\begin{proof}
Modulo $I_4$, the first two formal-sum coefficients are unit multiples of
$v_4,v_5$.  There is no decomposable correction involving $v_4$ to the
second of these coefficients: $u_5=pu_4+1$ is not a multiple of $u_4$.
The formal group law has no nonlinear term of total ordinary degree less
than $p^4$, because a coefficient of such a term would have positive weight
less than $u_4$.  Thus, through exponent $p^5$, its $p$-series is the ordinary
sum of its $x^{p^4}$ and $x^{p^5}$ terms.  The nonlinear terms in adding
these two inputs have exponent greater than $p^5$.

In the identity intertwining this $p$-series with a strict coordinate
change, only $x+t_1x^p$ from that change can affect these two exponents.
On the side where the change is applied after the $p$-series, the
coefficient of $x^{p^5}$ acquires a term proportional to $v_4^p t_1$.
On the other side it acquires a term proportional to
$v_4t_1^{p^4}$.  All nonlinear corrections to the strict formal-sum change
enter too late after taking the $p^4$-th power.  Consequently
\[
 \eta_R(v_5)-v_5=c\,v_4t_1^{p^4}+d\,v_4^p t_1
 \quad\text{in }\Gamma/I_4\Gamma,
\]
where $c,d\in\mathbb F_p$ and $c\ne0$.  The nonzero coefficient follows
also directly from the unit factors $c_4,c_5$ in
Equation~\eqref{fw:coordinate-unit}; their reductions cannot annihilate
the first displayed contribution.  Divide this coboundary by $c$ and put
$e=d/c$.  Since $t_1$ is primitive modulo $I_4$, all its $p$-power iterates
are primitive, and the resulting cohomology relation is
Equation~\eqref{fw:v5-cohomology-relation}.
\end{proof}

For a comodule concentrated in weights that are nonnegative integers, let
$C^{s,w}$ denote normalized cobar cochains of length $s$ and weight $w$.
We use the convention in which a positive-length cochain has a coefficient
followed by $s$ reduced cooperation factors.  This convention fixes the
signs of the usual cobar differential and cup product.

\begin{lemma}[Normalized positivity]
\label{fw:normalized-bound}
If the coefficient comodule has no negative weights, then
\begin{equation}
 C^{s,w}\ne0\quad\Longrightarrow\quad w\geq s.
 \label{fw:positivity}
\end{equation}
\end{lemma}

\begin{proof}
Every reduced cooperation factor is a polynomial monomial containing at
least one $t_i$, and hence has weight at least one.  The coefficient has
nonnegative weight.  Adding these weights proves the claim.
\end{proof}

\begin{lemma}[A comparison in three cochain degrees]
\label{fw:range-comparison}
Suppose $s\geq1$, $w<u_{m+1}$, and
$w-u_m<s-1$.  Additive specialization identifies the normalized cobar
complexes for $R/I_m$ and $P$, in weight $w$, in all three degrees
$s-1,s,s+1$.  It intertwines both adjacent differentials, and therefore
induces an isomorphism on cohomology in degree $(s,w)$.
The same conclusion in every cochain degree holds if $w<u_m$.
\end{lemma}

\begin{proof}
In each of the three degrees, any positive coefficient monomial consumes
at least $u_m$ weight.  It would leave at most $w-u_m<s-1$ for at least
$s-1$ reduced bar factors, which is impossible.  A coordinate $t_i$ with
$i\geq m+1$ exceeds the entire weight budget.  Thus these cochain groups
have scalar coefficients and only $t_1,\ldots,t_m$ in their bars.  Their
coproducts are additive by Lemma~\ref{fw:coordinates}.  This identifies
the two groups adjacent to the desired cohomology group as well as the
group itself, so both kernels and images agree.  If $w<u_m$, the same
exclusion of coefficients and coordinates works without a cochain-length
restriction.
\end{proof}

\subsection{Convergence, pairings, and the edge}

The following two propositions specify the topological spectral sequence and
the binary pairings used throughout the construction. The ordinary
Adams--Novikov resolution and its convergence are treated in
\citet[Theorems~2.2.3, 2.2.13--2.2.14]{Ravenel2004};
\citet[Theorem~2.3.3]{Ravenel2004} gives the binary-pairing, Leibniz,
and homotopy-compatibility statements in this form. Their complete proofs
in the particular resolution used here,
including the ordinary totalization model and its filtration, are given in
Appendix~\ref{app:resolution}.

\begin{proposition}[The bounded-below Adams--Novikov resolution]
\label{fw:anss}
For every bounded-below $p$-local spectrum $Z$, the sphere-source
Adams--Novikov spectral sequence has
\[
 E_2^{s,t}(Z)=H^{s,t}\bigl(C^\bullet_\Gamma(BP_*Z)\bigr)
 \quad\Longrightarrow\quad \pi_{t-s}Z.
\]
It converges strongly, and the induced filtration in each stem is finite.
The differential $d_r$ changes $(s,t)$ by $(r,r-1)$.  The filtration-zero
edge identifies $E_\infty^{0,*}$ with the image of the $BP$-Hurewicz map
$\pi_*Z\longrightarrow BP_*Z$.
\end{proposition}
The proof is given in Appendix~\ref{app:convergence}.


\begin{proposition}[Binary pairings of the resolution]
\label{fw:anss-pairing}
For bounded-below $p$-local spectra $Z_1,Z_2,Z_3$, an actual map
$f:Z_1\wedge Z_2\to Z_3$ induces pairings of the
Adams--Novikov spectral sequences, with the Leibniz rule on every page.
The $E_2$ pairing is the cobar cup pairing induced by the exterior
$BP$-homology pairing, and the $E_\infty$ pairing is the associated-graded
pairing of the actual homotopy map $f$.  No associativity of $f$ is needed.
If $Z_1=Z_2=Z_3=Z$ has a binary two-sided unital multiplication and
$BP_*Z=R/I$ is cyclic on its unit, the $E_2$ pairing is the ordinary
multiplication on cobar cohomology with quotient-ring coefficients.
\end{proposition}
The proof is given in Appendix~\ref{app:pairing}.


\begin{corollary}[Sparsity]
\label{fw:sparsity}
If $BP_*Z$ is concentrated in internal degrees divisible by $q$, then all
pages of its Adams--Novikov spectral sequence have that same internal
sparsity.  Every possibly nonzero positive differential has length
$r=qk+1$ for an integer $k\geq1$.  In weight notation it changes
$(s,w)$ to $(s+qk+1,w+k)$.  In particular, a primitive coefficient of
weight $u_n$ in filtration zero can have only targets
\begin{equation}
 (s,w)=(qk+1,u_n+k),\qquad k\geq1.
 \label{fw:primitive-targets}
\end{equation}
\end{corollary}

\begin{proof}
The coefficient and cooperation gradings give the assertion on the cobar
page.  A differential preserves the sparse support only if $q$ divides
$r-1$.  Subsequent pages are subquotients of earlier pages.  Since $r\geq2$
on or after the cohomology page, the first allowed length is $q+1$.
\end{proof}

\begin{lemma}[The additive comparison for the $v_4$ targets]
\label{fw:additive-comparison}
Let $s=qk+1$, $w=u_4+k$, $k\geq1$.  A nonzero normalized cobar group
in degree $(s,w)$ with coefficients $R/I_4$ requires
\begin{equation}
 (q-1)k\leq u_4-1,\qquad
 w\leq\frac{qu_4-1}{q-1}<2u_4<p^4<u_5.
 \label{fw:preliminary-v4-bound}
\end{equation}
For every such $k$, additive specialization is an isomorphism on
cohomology in degree $(s,w)$, by an identification of the cochain groups
of lengths $s-1,s,s+1$.  No $v_4$ coefficient or $t_4$ bar occurs in
any of those three cochain groups.
\end{lemma}

\begin{proof}
The first two inequalities are the rearrangement of $w\geq s$.
Here
\[
 u_4=1{,}028{,}262{,}820,\qquad
 2u_4=2{,}056{,}525{,}640<p^4=1{,}036{,}488{,}922{,}561,
\]
so the remaining inequalities in
Equation~\eqref{fw:preliminary-v4-bound} hold at the fixed prime.
For a cochain length $d\in\{s-1,s,s+1\}$, a coefficient containing
$v_4$ would leave at most $k$ weight for $d\geq qk>k$ reduced factors.
It therefore cannot occur.  A bar containing $t_4$ would leave at most
$k$ weight for the other $d-1\geq qk-1>k$ bars, and likewise cannot occur.
Higher generators exceed the weight bound.  The remaining coproducts
are additive by Lemma~\ref{fw:coordinates}; equivalently, apply
Lemma~\ref{fw:range-comparison}.  This identifies both adjacent
differentials and proves the cohomology assertion.
\end{proof}
